\subsection{\texorpdfstring{测度 \(\mu\) 关于 \(\mu\) 的伪像的分解}{测度 mu 关于 mu 的伪像的分解}}

定理 2. --- 设 T 与 B 是两个具有可数基的局部紧空间， \(\mu\) 是 T 上的正测度， \(p\) 是 T 到 B 中的 \(\mu\) -可测映射， \(\nu\) 是 \(\mu\) 在 \(p\) 下的一个伪像测度。则存在 T 上正测度的一个 \(\nu\) -适宜族 \(b \mapsto \lambda_b\) （ \(b \in \mathbf{B}\) ），它具有下列性质：

\begin{enumerate}
\def\labelenumi{\alph{enumi})}
\item
  \(\lambda_{b} \neq 0\) （对 \(b \in p(\mathrm{T})\) ）；
\item
  \(\lambda_{b}\) 集中于集合 \(\overline{p}^{-1}(b)\) 上（对每个 \(b\in \mathbf{B}\) ）；特别地， \(\lambda_{b} = 0\) （当 \(b\notin p(\mathrm{T})\) 时）；
\item
  \(\mu = \int \lambda_b d\nu(b)\)
\end{enumerate}

此外，若 \(\nu' = r \cdot \nu\) 是 \(\mu\) 在 \(p\) 下的另一个伪像测度，且 \(b \mapsto \lambda_b'\) 是 T 上正测度的一个 \(\nu'\) -适宜族，它关于 \(\nu'\) 具有性质 b) 与 c)，则 \(\lambda_b = r(b)\lambda_b'\) 在 B 中几乎处处成立（对 \(\nu\) 或 \(\nu'\) ）。

存在一个连续且有限的数值函数 \(f\) ，定义在 \(T\) 上，使得 \(f(t) > 0\) （对每个 \(t \in T\) ），且使得 \(\mu'' = f \cdot \mu\) 有界（Ch. V, §5, No.~6, 命题 11）。设 \(\nu'' = p(\mu'')\) ，它等价于 \(\nu\) ，记 \(\nu'' = g \cdot \nu\) ，其中 \(g\) 有限且局部 \(\nu\) -可积；此外还可以假定 \(g(b) > 0\) （对一切 \(b \in B\) ）（Ch. V, §5, No.~6, 命题 10）。把 No.~1 的定理 1 应用于 \(\mu''\) 与 \(\nu''\) ，就表明存在一个 \(\nu''\) -适宜族 \(b \mapsto \lambda_b''\) （ \(b \in B\) ）（ \(T\) 上正测度的族），使得：1) \(\| \lambda_b'' \| = 1\) （对 \(b \in p(T)\) ）；2) \(\lambda_b''\) 集中于 \(\overline{p}^{-1}(b)\) 上（对每个 \(b \in B\) ）；3) \(\mu'' = \int \lambda_b'' d\nu''(b)\) 。

对每个 \(b \in B\) ，我们定义一个正测度 \(\lambda_{b}\) （ T 上的），公式为 \(\lambda_{b} = (1/f) \cdot (g(b)\lambda_{b}^{\prime\prime})\) 。显然族 \(b \mapsto \lambda_{b}\) 具有叙述中的性质 a) 与 b)。另一方面，对每个函数 \(h \in \mathcal{K}(T)\) ， h/f 属于 \(\mathcal{K}(T)\) ，因此

\[
\int h (t) d \mu (t) = \int \left(h (t) / f (t)\right) d \mu^ {\prime \prime} (t) = \int d \nu^ {\prime \prime} (b) \int \left(h (t) / f (t)\right) d \lambda_ {b} ^ {\prime \prime} (t).
\]

但由于函数 \(b \mapsto \int (h(t) / f(t)) d\lambda_b''(t)\) 是 \(\nu''\) -可积的，函数 \(b \mapsto g(b) \int (h(t) / f(t)) d\lambda_b''(t)\) 是 \(\nu\) -可积的（Ch. V, §5, No.~3, 定理 1）。

根据 \(\lambda_b\) 的定义，这个函数就是 \(b \mapsto \int h(t) d\lambda_b(t)\) ，由此（出处同上）得 \(\int h(t) d\mu(t) = \int d\nu(b) \int h(t) d\lambda_b(t)\) ，这就证明了 \(\mu = \int \lambda_b d\nu(b)\) 。

为证明叙述的第二部分，我们注意可以假定 \(r(b) > 0\) （对一切 \(b \in B\) ）（Ch. V, §5, No.~6, 命题 10）；令 \(\lambda_b^{\prime \prime \prime} = f \cdot \left( \left( r(b) / g(b) \right) \lambda_b' \right)\) ；与上面一样可以证明，对每个函数 \(h \in \mathcal{K}(T)\) ，关系式

\[
\int h (t) d \mu (t) = \int d \nu^ {\prime} (b) \int h (t) d \lambda_ {b} ^ {\prime} (t)
\]

蕴含

\[
\int h (t) d \mu (t) = \int d \nu^ {\prime \prime} (b) \int \left(h (t) / f (t)\right) d \lambda_ {b} ^ {\prime \prime \prime} (t).
\]

因此，把 No.~1 的定理 1 应用于 \(\mu''\) 与 \(\nu''\) ，便推出对几乎每个 \(b \in B\) ， \(\lambda_b''' = \lambda_b''\) ，从而 \(\lambda_b = r(b)\lambda_b'\) 。

定义 2. --- 设 T 与 B 是两个 Polish 局部紧空间。给定 T 上的正测度 \(\mu\) 、 T 到 B 中的 \(\mu\) -可测映射 \(p\) 以及一个伪像测度 \(\nu\) （ \(\mu\) 在 \(p\) 下的伪像测度）， T 上正测度的每一个 \(\nu\) -适宜族 \(b \mapsto \lambda_b\) （ \(b \in B\) ），只要具有定理 2 的性质 \(b\) ) 与 \(c\) )，就称为 \(\mu\) 关于 \(\nu\) 的一个分解。

当 p 是 \(\mu\) -逆紧的且 \(\nu = p(\mu)\) 时，关于 p 的分解这一概念便与关于 \(\nu\) 的分解这一概念一致。在定理 2 的假设下， \(\mu\) 关于同一伪像测度 \(\nu\) 的两个分解关于 \(\nu\) 几乎处处相等。

注记. --- Ch. V, §3, No.~4 的定理 1 表明（考虑到 T 与 B 具有可数基这一事实）：对每个定义在 T 上、取值于 \(\overline{R}\) 或 Banach 空间 F 中且 \(\mu\) -可积的函数 f ，那些 \(b \in B\) 中使 f 不为 \(\lambda_{b}\) -可积者所成的集是 \(\nu\) -可忽略的，几乎处处有定义的函数 \(b \mapsto \int \mathbf{f}(t) d\lambda_{b}(t)\) 是 \(\nu\) -可积的，且

\[
\int \mathbf {f} (t) d \mu (t) = \int d \nu (b) \int \mathbf {f} (t) d \lambda_ {b} (t).
\]

对标量 \(\mu\) -可积的函数，应用 §1, No.~1 的命题 3，也有类似的结果。

